\documentclass[11pt,twocolumn,a4paper]{article}
\usepackage[utf8]{inputenc}
\usepackage[T1]{fontenc}
\usepackage{lmodern}
\usepackage{geometry}
\geometry{margin=1in}
\usepackage{graphicx}
\usepackage{booktabs}
\usepackage{amsmath,amssymb}
\usepackage{hyperref}
\usepackage{comment}
\usepackage[backend=biber,style=authoryear]{biblatex}
\addbibresource{references.bib}

% Title and Author
\title{At What Cost? A Comprehensive Transmission Costing Methodology for Informed Infrastructure Decisions}
\author{
  Andrew Igdal\thanks{University of Texas at Austin, LBJ School of Public Affairs} \\
}
\date{\today}

\begin{document}
\maketitle

% Abstract
\begin{abstract}
We introduce a comprehensive methodology designed to estimate and compare the costs of transmission infrastructure projects consistently and transparently.
% Provide a brief summary of the problem, the CTCC solution, key methods, and major findings.
\end{abstract}

% Section 1: Introduction
\section{Introduction}
% 1--2 pages. Motivation: U.S. clean electricity goals, grid constraints, cost overruns.

The U.S. electric grid is a vast, interconnected network of generators and loads, linked by extensive transmission and distribution (T\&D) systems. The grid has successfully electrified American households, businesses, and essential services, delivering substantial social and economic benefits. Central to this achievement are more than 160,000 miles of high voltage (69kV and above) transmission lines responsible for transporting electricity across vast distances \parencite{EIA2021Interconnections, OSHATransmissionDistribution2025}. However, today's transmission system is reaching its operational limits and in some locations, it already has. A rapid and significant expansion of transmission infrastructure is urgently needed over the next few decades to maintain grid reliability, ensure affordability, and meet ambitious  decarbonization targets.

A reliable electric grid meets electric demand and is able to withstand reasonable levels of disturbances  \parencite{ReliabilityExplainerFederal}. As a major - and the most exposed - part of the grid, adequate transmission infrastructure is essential for maintaining reliability. Power outages increased by 60\% between 2015 and 2020, resulting in over \$150 billion in costs to the U.S. economy \parencite{stoneCompoundClimateInfrastructure2021,CleanCheapterStrongerPEW}. Expanding transmission capacity is a critical step to reduce outages and enhance reliability \parencite{lawsonVariableRenewableEnergy} as the demand for power may double from 2020 to 2050\parencite{NationalTransmissionNeeds2023}. In events of reduced resource adequacy, an additional 35 GW of interegional transmission transfer capacity may be required by 2033 to secure grid reliability  \parencite{nercITCS2024}. In addition, it is important for electricity to remain affordable.

An affordable grid delivers electricity at the lowest possible cost to consumers. Transmission investments can reduce overall costs through several key mechanisms. First, an expansion of transmission capacity reduces congestion—constraints. Congestion occurs when existing transmission capacity is at risk of being overloaded, making it unable to carry more power. The resulting congestion prevents the dispatch of cheaper power, which means more expensive generators are required to be dispatched to match local supply and demand - increasing prices. With greater transmission capacity, congestion can be reduced or alleviated, decreasing prices. Second, transmission infrastructure connects geographically distant but economically advantageous sources of electricity, such as wind and solar resources with demand centers that might be concentrated elsewhere. Third, by expanding transmission capacity, the grid reduces curtailment - intentional reductions in wind and solar generation  due to local oversupply — and maximizes the utilization of low-cost generation. The economic benefits of transmission expansion can significantly outweigh associated costs. For example, facilitating the integration of affordable wind and solar generation could reduce a typical U.S. household electricity bill by over \$25 per month or approximately \$300 per year, representing about 3.6\% of residential consumer electricity costs \parencite{clackConsumerEmploymentEnvironmental}. In aggregate, each dollar spent on transmission can result in savings ranging from \$1.60 to \$1.80, translating into potential cumulative savings approaching \$1 trillion in electricity system cost savings by 2035 \parencite{DOE2024NTPS}.

Transmission expansion also directly supports state, regional, and national decarbonization goals.  The necessary scale and pace of transmission expansion must accelerate significantly compared to historical trends to meet national decarbonization objectives \parencite{DOE2024NTPS}.

Maybe something along the lines as Tx acting as an economic growth facilitator? Case: CREZ unlocked wind, which then electrified O\&G and now TX is building more Tx (Permian Reliability Plan) to unlock more O\&G, which will unlock more solar. RE and O\&G devs couldn't do it themselves, so the state decided to develop the conditions through Tx that will make therm successful?

% Gap: Existing tools lack comprehensive cost accounting and scenario flexibility.

Despite evidence supporting expanded transmission, actual progress has been slow. Construction of high-voltage transmission lines during the early 2020s has decreased by 20\% compared to the first half of the 2010s. From 2010 to 2014 the average miles of transmission built per year was 1,700, that has fallen to 350 miles per year between 2020-2023 \parencite{FEWERMILES2024}. There are many explanations and barriers to transmission expansion. Among the barriers is so called "not in my back yard" public opposition (NIMBYism), ineffective planning and inadequate coordination, and uncertainty regarding projected costs and benefits (that can often make benefit to cost ratios of projects lower than they might be in reality) \parencite{joskowFacilitatingTransmissionExpansion2022}. This uncertainty is exacerbated by the lack of standardized methodologies for comprehensively valuing transmission projects, which frequently leads to project delays, cost overruns, and cancellations.

% Introduce CTCC: Side-by-side, nominal vs. real, transparent constants.
% Objectives and research questions.

This paper introduces a flexible and comprehensive methodology for scenario-based transmission cost accounting. It addresses the existing knowledge gap in flexible, standardized methods by presenting a detailed, consistent, and rigorous generalized framework for evaluating transmission projects under multiple technologies, terrain types, risk factors, financing options, permitting regimes, development timelines, and routing options. We demonstrate the practical application of this framework and discuss future development objectives for the Comprehensive Transmission Costing Calculator (CTCC). The remainder of the paper is structured as follows: Section 2 provides background on existing cost estimation approaches, including their strengths and their limitations. Section 3 details the methodological framework and structure of the CTCC. Section 4 presents illustrative case studies comparing different transmission project scenarios. Section 5 discusses implications, potential applications, and current limitations of the CTCC. Section 6 concludes the paper by summarizing key contributions and outlining future research directions.

% Section 2: Background
\section{Background}

% 2--5 pages. 

% Discuss hard vs. soft costs, avoided costs, and decision-making frameworks.
% Review literature on transmission cost estimation.

\subsection{Introduction to Transmission Costing}
When decision makers are planning transmission projects, approval often hinges on cost. Thus, accurate costing of transmission projects is crucial. The costs of a project can be divided into three main categories: hard costs, soft costs, and avoided costs. Hard costs are tangible expenses directly related to the construction and operation (materials, wages, engineering, land, etc.) and are often the easiest to track. Soft costs, unlike hard costs, are indirect costs and often intangible until paid for (permitting, mitigation, legal fees, administrative overhead, financing). Avoided costs, like soft costs, are often indirect and intangible. Costs not incurred due to investments in transmission like reduced outages, decreased wildfire risks, avoided curtailments, and congestion reductions, are all examples of avoided costs. Estimating these costs can be difficult and are often ignored \parencite{EllramTate2021}. Most methods for estimating transmission costs focus on hard costs and soft costs; few methods consider avoided costs - no method standardizes consideration for all three.

\subsection{Factors Influencing Transmission Cost Variability}
There is wide variability in transmission project costs. For example, Texas' Competitive Renewable Energy Zone (CREZ) transmission lines had an average cost of around \$2,500 per MW-mile, while California's Sunrise Powerlink had an average cost of around \$16,000 per MW-mile \parencite{AndradeBaldick2017}. The variation in these costs depends on various factors. First, the terrain and geography of any particular project greatly influences costs. Construction costs in mountainous, forested, or wetland terrain is generally more expensive than flat and open land. MISO's Transmission Cost Estimation Guide estimates that forest clearing costs \$6K/acre, independent of the biome or topography; building on wetland areas (often forested) adds an additional \$70K/acre for matting, and construction difficulties on top of compensatory wetland mitigation credits (i.e., fees paid to offset wetland impacts) under the Clean Water Act costing \$56K/acre. Mountainous terrain may add additional costs of \$8K/acre. \parencite{MISO2024CostGuide}. Second, the level of development on land needed for a project under consideration affects cost. Installing lines near cities or other developed lands often costs more acre. Acquiring right of way (ROW) in urban areas can be costly and controversial. Although most of Texas' CREZ lines stretch for hundreds of miles along open, undeveloped land, some lines were necessarily constructed in urban areas, leading to increased costs and delays \parencite{AndradeBaldick2017}. Third, the acquisition of ROW can contributes massively to costs and delays. Acquiring ROW involves negotiations, appraisals, counter-appraisals, and sometimes eminent domain, often from multiple landowners and interested parties \parencite{IUB2021EminentDomainFAQ}. Costs can vary widely per acre, from a few hundred dollars per acre of pasture land to tens of thousands of dollars per acre of urban land \parencite{MISO2024CostGuide}. These costs can sum to millions per mile. Utilities often prefer to avoid condemnation via eminent domain \parencite{WI_DATCP_ROW}. The development of the Grain Belt Express (an 800-mile High Voltage Direct Current (HVDC) line) shows that relying on eminent domain can often lead to court battles and costly delays \parencite{Vanacker2021GrainBeltCondemnation}. 

Fourth,  labor and material costs - while often relatively easier to calculate - can vary widely from project to project and even over time for the same project. For example, the Wisconsin Cardinal-Hickory Creek project saw double digit percent increases in conductor and labor costs, and triple digit percent increases in steel pole prices from its first estimate in 2018 \parencite{Kaeding2023CardinalHickory}. Fifth, regulatory and permitting complexities can cause lengthy delays and inflate project costs. Many projects face legal challenges during lengthy federal/state processes that increase legal costs and force planners to continuously revise routes and options. One study found that substantial opposition contributed to delays and cancellations in ~27\% of large U.S. transmission proposals studied \parencite{BaranoffNorris2024}. For projects that are not canceled due to these hurdles, costs often rise during these delays. Sixth, the overall scale and complexity of a project is a large driver of cost in both higher and smaller directions. Larger projects may benefit from economies of scale, or their sheer scale may contribute to extra costs and/or delays that increase costs. Using new technologies might have a similar result. For example, Texas' CREZ lines' large scale led to some cost savings per MW-mile. Yet, CREZ's final cost of \$6.9B was 40\% higher than its originally estimated cost of \$4.9B. Isolated lines like the Sunrise-Powerlink are unable to take advantage of scale and face significant costs due to complexity associated with high capacity lines and its chosen route (i.e. crossing mountains and undergrounding in sensitive areas, etc.) \parencite{AndradeBaldick2017}. 


Taken together, there are many uncertainties regarding cost estimation: terrain and the environment, regulation, location, labor and materials, ROW, and project complexity create circumstances that can result in lengthy and costly delays. 

\subsection{Existing Tools and Their Limitations}
Planners need tools that combine each of these costs, with an ability to carefully calibrate to their projects' specific conditions to estimate costs. Although useful tools for calculating transmission costs exist in both academia and industry,  we found none captures a full range of costs. The following section describes current tools and methods for costing transmission, their strengths and weaknesses, and highlights contributions that the CTCC provides to the field of transmission cost estimation tools and models.


\subsection{Academic Tools}
In the academic literature, multiple methods exist for estimating the costs and benefits of transmission infrastructure. Examples include Larsen's novel attempt to standardize a method for estimating the costs and benefits of undegrounding transmission, as well as more complex (and harder to utilize) methods such as Teegala and Singal's work to optimally estimate hard costs of transmission via genetic algorithms (\parencite{Larsen2016, TeegalaSingal2016}. These methods are highly specific and less generalizable than is needed for specific project planning.

Another method derived from academic and lab work for estimating the costs of grid expansion (including generation, load, and transmissions) are capacity-expansion models. These models are typically used in national labs and universities. Examples include the National Renewable Energy Laboratory's Regional Energy Deployment System (ReEDS) and MITs GenX \parencite{NREL_ReEDS_2019Doc, JenkinsSepulveda2017_GenX}. Both models are able to expand grid infrastructure from an initial state to a final state based on a future load profile and possible measures. Included in these sorts of simulations are the actual transmission needs. As such, transmission is given some estimated cost per MW-Mile. Such estimates vary largely in theory and in practice. When using such capacity expansion models, practitioners must exogenously choose such an \$ per MW-mile estimate for the model. This provides a consistent high-level estimation of costs and needs for transmission expansion. However, it lacks granularity as there is often little to no variation in land costs, construction costs, permitting times, outage risks taken into account.  Unlike GENX, which uses one MW-mile cost assumption over an entire simulation, ReEDs utilizes a more complex estimation for transmission costs by creating a high-geographic-resolution cost surface \parencite{LopezEtAl2025_NREL_TechnicalPotential}. However, it is still missing many other costs and benefits that play a role . Overall, these models are designed for high-level regional or national scenario analysis, and are useful for broad, strategic planning, but insufficient for detailed, project-level decision making. Additionally, their technical nature often make them inaccessible and complicated.

The Reconductoring Economic and Financial Analysis (REFA) tool developed by Lawrence Berkely National Laboratory (LBNL), is a step toward improving the estimation and comparison of transmission project costs and benefits. REFA provides a focused cost-benefit analysis of reconductoring projects. It considers the full net present value of a project and the benefits of reduced line losses, and can help demonstrate long-term justifications for investments. Though more intuitive and accessible, the REFA tool has inherent limitations given its purpose. It targets reconductoring scenarios, addressing hard cost components - not including the range of soft costs that may be present \parencite{LBNL_REFA_Manual_2025}.


\subsection{Industry Tools}
Industry has also worked to develop methods to estimate the cost of transmission. Black and Veatch's transmission infrastructure capital cost tool was created to estimate costs of projects throughout the Western Electric Coordinating Council (WECC) region (Pletka Paper). This tool is extensive and was used as a basis for the creation of the CTCC. It considers capital costs associated with conductor type, strucutre type, substations, inflation, conductor type, structure type, line lengths, varying terrains, and varying right-of-way (ROW) sizing. It also allows for the consideration of reconductoring, replacing wires on a circuit, rather than just new builds. Black and Veatch's method improved the reporting granularity of hard costs associated with transmission infrastructure as well as transparency and accessibility, but lacks any costing of soft costs. It also does not consider any avoided costs. Another available tool is PLS-CADD. Its primary use is for detailed line design and cost estimation. It includes detailed databases of materials, parts, and assemblies allowing for detailed hard costing of a project \parencite{PLSCADD2012}. It also allows for the custom addition of fields such as engineering labor, ROW costs, and other construction related expenses. However, it does not estimate comprehensive overhead or soft costs or benefits. PLS-CADD excels at detailed structural and material cost estimation but is limited in comprehensive scenario comparison.

Independent Systems Operators (ISOs) and Regional Transmission Operators (RTOs) are organizations tasked with running the grid, maintaining resource adequacy, and planning for the future. These organizations have developed methods for estimating the costs of transmission, and these methods vary from entity to entity. The Midcontinent ISO (MISO) has a transparent cost modeling framework detailed in its Transmission Cost Estimation Guide. They break costs into components such as project management costs, administrative and general overhead, engineering and environmental studies, ROW and land acquisition, regulatory and permitting costs, as well as the hard costs associated with the construction of a transmission project \parencite{MISO2024CostGuide}. While their methods include more soft costs than previous methods, they make fixed assumptions for the proportion of costs the soft costs will make up. They also use a reasonable, but generally inflexible, accounting method for ROW and land acquisition costs. MISO's method remains mostly an engineering estimate template that does not take into account project-specific factors, risks, and potential benefits.(CITE)

PJM, another RTO, relies on competitive proposals from developers, such that cost estimates come in as ranges (Sliwa power point). Thus, there is no single PJM method. Instead, PJM can observe the bids and their range, and compare it to historical data to judge for realism. This process has benefits; price variability can be revealed (proposals can range from a few million to several billion). By relying on historic data and market bids, PJM can effectively capture information other more structured and less flexible methods cannot \parencite{Tesfa2024PJM}. Still, without a consistent model, PJM cannot independently explore factors other than construction. Without a shared common model between ISOs and RTOs, planning and cost estimation has generally remained fragmented and inconsistent. Reliance on rule-of-thumb estimates of \$ per MW-mile or case-by-case studies often don't incorporate the whole breadth of costs. 

\subsection{The Need for Comprehensive Tool}
Current methods are useful but can be lacking in four key areas. First, a lack in scenario flexibility - highlighted models are often rigid and fixed, and those with relatively more flexibility, such as ReEDs, Genx, Black and Veatch, and MISO still rely on inflexible assumptions, and are too focused on hard costs. Second, there is a lack of soft cost consideration. Soft costs are often missing, overly simplified through fixed multipliers. Soft costs have high variability and are often extremely sensitive to delays - none of the methods we found consider the massive increase in soft costs delays can incur. Third, avoided costs (or benefits) are generally under-considered. When considered, most models focus on decreases in curtailment and congestion but little else. Fourth, and most importantly, while there do exists tools that consider soft costs, or avoided costs, or allow for some flexibility in scenarios, no one tool combines all into a simple, accessible, and standardized method. This fragmented tool environment makes cross-project comparison more difficult, and can complicate planning and investment decisions. 

The importance of accurate, standardized and comprehensive cost estimates cannot be overstated. When project costs are under-projected, budgets may not support the project and ratepayers and taxpayers will ultimately bear unexpected overruns. On the other hand, over-projected costs are also harmful as they can deter investment from ever occurring. FERC's order 1920 highlights concern from the highest level of uncertain costing of transmission projects. The order requires any regional transmission projects be reevaluated if its actual or project cost "significantly exceed[s] the cost estimates used during selection" \parencite{FERC2024JointConcurrence1920}. If a project’s costs balloon beyond initial projections, it may no longer be the most efficient or cost-effective solution and might not meet project selection criteria. Thus, accurate upfront estimates of costs and benefits are important to algin transmission investment with transmission needs. 

\subsection{CTCC Contribution and Innovation}
In summary, existing transmission costing methods are fragmented and insufficiently comprehensive, often neglecting critical soft and avoided cost components and lacking flexibility. The CTCC addresses these gaps by integrating all relevant cost categories into a unified, flexible, and transparent framework. By enabling scenario comparisons with immediate feedback, the CTCC can improve decision-making efficiency, transparency, and cost estimation accuracy, directly benefiting planners, investors, regulators, and affected communities alike.





% Identify shortcomings in current methodologies and tools.



\section{Methods}
% 3--5 pages. Detail CTCC architecture:
% - Input validation and parameter derivation.
% - Cost modules: Build, ROW, delay, O&M, soft, emissions, avoided costs.
% - Nominal-to-real conversion (discounting).
% - Aggregation and side-by-side scenario analysis.
% - User-adjustable constants and transparency.
% Include equations where relevant (e.g., net present cost calculation).


The CTCC combines hard, soft, and avoided costs in one place. It allows for a near-infinite amount of scenarios to be tested side by side to facilitate quick and efficient planning and comparison of projects and project alternatives. The CTCC transparently and consistently facilitates the estimation of:

\begin{itemize}
  \item Different asset lifetimes, permitting, and construction timelines;
  \item Various construction and technology types;
  \item Construction costs;
  \item Right-of-way (ROW) costs;
  \item Terrain-specific costs;
  \item Congestion and curtailment reduction benefits;
  \item Avoided wildfire and outage costs;
  \item Delay costs;
  \item Insurance costs;
  \item Permitting and regulatory costs;
  \item Allowance for Funds Used During Construction (AFUDC);
  \item Social emissions costs associated with line losses.
\end{itemize}

By incorporating these components, the CTCC can provide both the total and net lifetime costs of transmission projects, presented in nominal and real terms. The following subsections provide a high-level overview of the methodology used by the CTCC. Detailed mathematical formulations for all cost categories, including nominal-to-real conversions, conditional logic for different project types, and user-adjustable parameters, are provided in Appendix A.

\subsection{Methodological Framework}

The CTCC framework aggregates costs and benefits into clearly defined categories that span the entire project lifecycle. The methodology distinguishes between three temporal phases: the permitting phase (duration $Y$), the construction phase (duration $T_{\mathrm{constr}}$), and the operational phase (duration $N$). All costs and benefits are calculated in both nominal terms and present value, using appropriate discount rates for financial and societal perspectives.

The total system cost in present value terms follows the general structure:
\begin{equation}
\text{Total\_Cost\_PV} = \text{Capital\_Costs\_PV} + \text{Operational\_Costs\_PV} + \text{Energy\_Emissions\_Costs\_PV} + \text{Risk\_Costs\_PV} + \text{Delay\_Costs\_PV} - \text{Benefits\_PV}
\label{eq:total_cost_framework}
\end{equation}

Each major category encompasses multiple subcomponents, as detailed in the following subsections. The framework is designed to handle conditional logic for different project types (e.g., greenfield vs. reconductoring, AC vs. DC, overhead vs. underground), with detailed formulations provided in Appendix A.

\subsection{Cost Categories}


\subsubsection{Capital Costs}

Capital costs represent upfront investments required to bring the transmission project into service. These include:

\textbf{Build Costs:} Construction costs for conductors, structures, and converter stations (for DC projects). Costs are calculated per mile and adjusted for terrain complexity using weighted miles. The framework distinguishes between greenfield projects and reconductoring projects, where reconductoring eliminates structure and converter costs. For DC projects, converter station costs are added. All build costs are subject to terrain multipliers that account for construction difficulty across different terrain types. See Appendix A.1 for detailed formulations.

\textbf{Right-of-Way (ROW) Costs:} Land acquisition, holding, and rental costs. For greenfield projects, ROW costs include one-time acquisition fees, annual holding costs during the permitting phase, and annual rental costs throughout construction and operation. For reconductoring projects, only rental costs apply since existing ROW is utilized. ROW costs vary by geographic zone and are calculated based on acreage requirements, which depend on line length and ROW width specifications. See Appendix A.2 for detailed formulations.

\textbf{Environmental Mitigation Costs:} Costs associated with environmental impact mitigation, including base mitigation costs that vary by construction type and terrain, wetland mitigation credits, and habitat mitigation credits. These costs are incurred during the construction phase and may be subject to AFUDC capitalization. See Appendix A.3 for detailed formulations.



\subsubsection{Operational Costs}

Operational costs represent recurring expenses throughout the project's operational lifetime:

\textbf{Operation \& Maintenance (O\&M):} Annual maintenance costs for conductors, structures (for overhead projects), converters (for DC projects), and vegetation management (for overhead projects). O\&M costs are calculated per mile and vary by component type and terrain characteristics. See Appendix A.4 for detailed formulations.

\textbf{Insurance Costs:} Annual insurance premiums on insurable assets (conductors, structures, and converters). Premium rates vary by construction type and are applied to the insured value of assets. See Appendix A.5 for detailed formulations.

\subsubsection{Energy and Emissions Costs}

\textbf{Line Loss Costs:} The economic cost of energy lost during transmission, calculated based on $I^2R$ losses for AC lines and resistance-based losses for DC lines. For reconductoring projects, line loss costs represent the difference between old and new configurations (typically a benefit). Losses are calculated considering line utilization, capacity, voltage, and operating hours. See Appendix A.6 for detailed formulations.

\textbf{Emissions Costs:} The societal cost of additional generation required to compensate for line losses, accounting for CO$_2$, SO$_x$, and NO$_x$ emissions. Emissions are calculated based on the energy mix of the grid, which may evolve over the project lifetime. The framework uses social discount rates for emissions costs to reflect societal preferences. See Appendix A.7 for detailed formulations.

\subsubsection{Risk Costs}

Risk costs capture expected annual losses from low-probability, high-impact events:

\textbf{Wildfire Costs:} Expected annual loss (EAL) from transmission line-caused wildfires, calculated as the product of ignition rates (which vary by terrain and construction type) and mean loss per event. The framework models wildfire risk as a growing annuity to account for increasing wildfire severity over time. See Appendix A.8 for detailed formulations.

\textbf{Outage Costs:} Expected annual cost (EAC) from transmission outages, calculated based on outage rates (which vary by construction type and terrain), outage duration, capacity at risk, and value of lost load (VoLL). VoLL is modeled using piecewise-linear functions that reflect higher costs for larger outages. See Appendix A.9 for detailed formulations.

\textbf{Wildfire Liability Insurance:} Optional annual premiums for wildfire liability coverage, applied when liability insurance is enabled. See Appendix A.10 for detailed formulations.

\subsubsection{Delay Costs}

Delay costs represent expenses incurred during the permitting phase, including legal processes, administrative overhead, contractor and labor costs, material and equipment costs, permitting and regulatory compliance, public relations, project management, environmental studies, and miscellaneous expenses. These costs are incurred annually during the delay period and are discounted to present value. See Appendix A.11 for detailed formulations.

\subsubsection{Capitalization Costs}

\textbf{Allowance for Funds Used During Construction (AFUDC):} The cost of capital during the construction period, calculated based on the weighted average cost of capital (WACC) and construction duration. AFUDC is applied to capital costs according to timing patterns that specify when costs are incurred during construction. See Appendix A.12 for detailed formulations.

\textbf{Contingency:} A percentage markup applied to build costs to account for cost overruns and uncertainty. Contingency rates are user-defined and may vary by component (conductor, structure, converter). See Appendix A.12 for detailed formulations.

\subsection{Benefits}

\subsubsection{Congestion Reduction Benefits}

Congestion reduction benefits capture the value of avoided transmission bottlenecks. The framework calculates annual congestion benefits as the product of energy congestion reduction (MWh/year) and average congestion price (\$/MWh). The methodology accounts for capacity allocation between curtailment and congestion to avoid double-counting, and includes saturation factors to reflect diminishing marginal benefits. The framework also calculates residual congestion (unrelieved congestion) and opportunity costs during delay periods. See Appendix A.13 for detailed formulations.

\subsubsection{Curtailment Reduction Benefits}

Curtailment reduction benefits capture the value of avoided renewable energy curtailment. Annual curtailment benefits are calculated as the product of curtailment reduction (MWh/year) and average curtailment price (\$/MWh). The framework includes saturation factors to account for diminishing marginal benefits as more capacity is added. See Appendix A.14 for detailed formulations.

\subsubsection{Revenue Requirements}

For rate-regulated utilities, the framework can calculate revenue requirements from a firm perspective. When rate-based revenue is enabled, the framework calculates annual revenue requirements as the product of the capital cost rate base and an allowed return rate. This represents the revenue stream needed to recover capital investments over the project lifetime. Revenue requirements are calculated in present value terms and are subtracted from total system costs when evaluating projects from a firm perspective. See Appendix A.17 for detailed formulations.

\subsection{Cost Aggregation}

The CTCC aggregates all cost categories into total lifetime costs and net lifetime costs. Total lifetime costs (TLC) represent the sum of all costs:
\begin{equation}
\begin{split}
\text{TLC} = \text{DC} + \text{ENM} + \text{ROW} + \text{Build} \\
+ \text{OM}            + \text{Insur} + \text{Cap} + \text{Emis} + \text{Risk}
\end{split}
\label{eq:TLC}
\end{equation}

where DC represents delay costs, ENM represents environmental mitigation costs, and Risk represents the sum of wildfire, outage, and wildfire liability insurance costs.

Net lifetime costs (NLC) subtract benefits:
\begin{equation}
\text{NLC} = \text{TLC}  - \text{B}
\label{eq:NLC}
\end{equation}
where B represents the benefits of congestion and curtailment reduction, as well as revenue to the firm.

All costs and benefits are calculated in both nominal terms and present value. Present value calculations use the weighted average cost of capital (WACC) for financial costs and the social discount rate for societal costs (emissions). The framework allows users to specify different discount rates for different cost categories, providing flexibility for different analytical perspectives.

\subsection{Scenario Flexibility and Transparency}

The CTCC framework is designed for maximum scenario flexibility. Users can specify different construction types (overhead, underground direct-buried, underground tunnel, subsea, reconductoring), technology types (AC vs. DC), conductor types, terrain compositions, financing structures, and risk parameters. All user-adjustable parameters are transparently documented, and the framework provides both regulatory perspectives (using AFUDC) and societal perspectives (using present value discounting).

Conditional logic throughout the framework ensures that costs are calculated appropriately for each scenario type. For example, reconductoring projects eliminate structure costs and ROW acquisition costs, while DC projects add converter costs and converter O\&M. The detailed conditional logic and all mathematical formulations are provided in Appendix A.



% Section 4: Analysis
\section{Analysis}
% 2--4 pages. Present an illustrative case study:
% - Example projects (overhead vs. underground vs. reconductoring).
% - Tables and figures comparing costs (nominal and present-value).
% - Discussion of trade-offs observed.

% Do something like figure 1 from Josh's paper:  https://www.sciencedirect.com/science/article/pii/S0301421516306875

\subsection{Analytical Framework}

To evaluate transmission project scenarios, we employ multiple analytical approaches that provide complementary perspectives on project economics and risk. This subsection outlines the benefit-cost ratio (BCR) metrics, sensitivity analysis methodology, and one-at-a-time (OAT) analysis approach used throughout this analysis.

\subsubsection{Benefit-Cost Ratio Metrics}

We calculate and present multiple BCR metrics to capture different analytical perspectives and decision-making contexts:

\textbf{BCR\textsubscript{system}:} The primary system-level metric, calculated as total benefits divided by total costs (both in present value). This comprehensive metric includes all costs (capital, operational, energy/emissions, risk, and delay) and all benefits (congestion reduction, curtailment reduction, revenue requirements, and line loss benefits for reconductoring projects).

\textbf{BCR\textsubscript{capital}:} Benefits divided by capital costs only, providing an investor-focused perspective that isolates the return on upfront investment.

\textbf{BCR\textsubscript{capital and delay}:} Benefits divided by capital costs plus delay costs, providing a more complete investor perspective that captures all pre-operational expenses (capital investment and permitting/regulatory delays). This metric addresses the gap where delay costs, while incurred during the pre-operational period, are not included in the capital costs category, offering a more comprehensive view of upfront investment requirements.

\textbf{BCR\textsubscript{excluding risk}:} Benefits divided by costs excluding risk components (wildfire, outage, and wildfire liability insurance). This metric focuses on core economic viability by removing low-probability, high-impact risk events that may be subject to different risk management strategies.

\textbf{BCR\textsubscript{excluding emissions}:} Benefits divided by costs excluding energy and emissions costs (line losses and emissions compensation). This metric isolates the direct project economics from environmental externalities.

\textbf{BCR\textsubscript{excluding emissions and risk}:} Benefits divided by costs excluding both energy/emissions costs and risk costs. This metric provides a core economic assessment by removing both environmental externalities and low-probability risk events, focusing on the fundamental project economics of capital, operational, and delay costs relative to system benefits.

All BCR metrics use conservative benefit estimates, with benefit values adjusted by saturation factors to account for diminishing marginal returns as transmission capacity is added. This provides a more conservative assessment of project benefits across all metrics.

All BCR metrics are calculated using present value terms, with benefits and costs discounted using appropriate rates (WACC for financial costs, social discount rate for emissions).

\subsubsection{Sensitivity Analysis Methodology}

To identify which parameters most significantly influence project economics, we employ Latin Hypercube Sampling (LHS) for systematic parameter variation. LHS ensures efficient coverage of the parameter space while maintaining statistical properties that enable robust sensitivity analysis.

For each scenario, we generate $N$ parameter samples using LHS, where each sample represents a complete set of parameter values drawn from their specified distributions. The CTCC is then evaluated for each sample, producing distributions of BCR metrics and cost/benefit components.

To quantify parameter importance, we calculate Partial Rank Correlation Coefficients (PRCC) between each input parameter and each output metric. PRCC measures the monotonic relationship between a parameter and an output while controlling for the effects of all other parameters, making it well-suited for non-linear models with correlated inputs. Parameters with high absolute PRCC values (typically $|\text{PRCC}| > 0.3$) are considered influential and warrant further investigation through OAT analysis.

\subsubsection{One-at-a-Time Analysis}

For the most influential parameters identified through sensitivity analysis, we conduct one-at-a-time (OAT) parameter sweeps. For each top parameter (typically the top 6 parameters per BCR metric, ranked by absolute PRCC), we systematically vary that parameter across a range of multiplier values while holding all other parameters at their baseline values.

This approach produces response curves showing how each BCR metric changes as a function of the parameter of interest. OAT analysis provides intuitive visualization of parameter impacts and helps identify critical thresholds or non-linear relationships that may not be apparent from PRCC values alone. The results are presented as plots showing BCR response curves for each influential parameter, enabling stakeholders to understand which parameters drive project economics and where uncertainty reduction efforts would be most valuable.


% Section 5: Discussion
\section{Discussion}
% 1--2 pages. Interpret results:
% - What are the implications for planners, regulators, and stakeholders given the scenarios and the results? Different options? .
% - How does the CTCC improve decision making for these scenarios? Are non-obvious decisions the better decision?.
% - Limitations of the current tool version - no amortization, .

% Section 6: Conclusion
\section{Conclusion}
% 1--2 pages. Summarize contributions:
% - Comprehensive cost framework.
% - Potential impact on transmission planning.
% Suggest future work: testing with stakeholders, integration with GIS, expanded cost modules.

















\end{document}
